% !TEX root = ../main.tex
\begin{ouverture}[acidebase]{1}{Transformations acide-base}{Partie 1 · Constitution et transformations de la matière}
{Acides et bases de Brönsted\enspace·\enspace Couples\enspace·\enspace Espèces amphotères\enspace·\enspace pH}
\objectif{Identifier un transfert d'ion hydrogène et les couples acide-base mis en jeu.}
\objectif{Établir l'équation d'une réaction acide-base à partir de deux demi-équations.}
\objectif{Représenter le schéma de Lewis d'un acide carboxylique, d'un ion carboxylate, d'une amine et d'un ion ammonium.}
\objectif{Reconnaître une espèce amphotère.}
\objectif{Passer du pH à la concentration en ion oxonium, et inversement.}
\prerequis{Schémas de Lewis, électronégativité et polarité d'une liaison.}
\prerequis{Quantité de matière, concentration et tableau d'avancement.}
\prerequis{Puissances de dix et notation scientifique.}
\end{ouverture}
\essentiel
\essSection{1}{Acides et bases selon Brönsted}
\begin{definition}[Acide, base, couple acide-base]
Un \textbf{acide} est une espèce chimique capable de céder un ion hydrogène \ce{H+}.
Une \textbf{base} est une espèce chimique capable de capter un ion hydrogène \ce{H+}.
Un acide et une base qui se transforment l'un en l'autre par perte ou gain d'un ion \ce{H+}
sont dits conjugués : ils forment un \textbf{couple acide-base}, noté acide/base, auquel on associe
une demi-équation :
\[\text{acide}=\text{base}+\ce{H+}.\]
\end{definition}
\begin{exemple}
\begin{tabular}{@{}l@{\hspace{6mm}}l@{}}
Couple \ce{CH3COOH}/\ce{CH3COO-} : & \ce{CH3COOH = CH3COO- + H+}\\[0.6mm]
Couple \ce{NH4+}/\ce{NH3} : & \ce{NH4+ = NH3 + H+}
\end{tabular}\par\vspace{1mm}
L'acide et la base d'un couple ne diffèrent que par un ion \ce{H+} : la charge de l'acide
dépasse d'une unité celle de sa base conjuguée.
\end{exemple}
\begin{attention}
Le signe = d'une demi-équation traduit une correspondance formelle. L'ion \ce{H+} n'existe pas
seul en solution aqueuse : il est toujours capté par une base, par exemple par l'eau,
qui forme alors l'ion oxonium \ce{H3O+}.
\end{attention}
\essSection{2}{Une réaction acide-base est un transfert d'ion hydrogène}
\begin{propriete}[Réaction acide-base]
Une réaction acide-base met en jeu deux couples : l'acide de l'un cède un ion \ce{H+}
à la base de l'autre.
\[\text{acide}_1+\text{base}_2\longrightarrow\text{base}_1+\text{acide}_2\]
\end{propriete}
\schema{01-transfert}
\begin{savoirfaire}[Établir l'équation d'une réaction acide-base]
\begin{questions}
\item Repérer l'acide, qui perd un ion \ce{H+}, et la base, qui en gagne un ; écrire les deux couples.
\item Écrire la demi-équation de l'acide dans le sens où il cède \ce{H+}, celle de la base
dans le sens où elle capte \ce{H+}.
\item Additionner les deux demi-équations : les ions \ce{H+} s'éliminent. Vérifier la conservation
des éléments et de la charge.
\end{questions}
\[\begin{array}{l}
\ce{CH3COOH = CH3COO- + H+}\\
\ce{NH3 + H+ = NH4+}\\\hline
\ce{CH3COOH + NH3 -> CH3COO- + NH4+}
\end{array}\]
\end{savoirfaire}
\begin{propriete}[Reconnaître un transfert d'ion hydrogène]
Une variation du pH, le changement de teinte d'un indicateur coloré ou un dégagement de dioxyde
de carbone (lorsque les ions hydrogénocarbonate ou carbonate captent des ions \ce{H+}) signalent
souvent une réaction acide-base. La flèche $\longrightarrow$ est réservée aux transformations
totales ; le symbole $\rightleftharpoons$ indique une transformation non totale, au cours de laquelle
la réaction se produit dans les deux sens (chapitre 7).
\end{propriete}
\essSection{3}{Couples à connaître et espèces amphotères}
\begin{center}\renewcommand{\arraystretch}{1.3}\small
\begin{tabularx}{\linewidth}{@{}lYY@{}}\toprule
\otTete{Couple} & \otTete{Acide} & \otTete{Base}\\\midrule
\ce{H3O+}/\ce{H2O} & ion oxonium & eau\\
\ce{H2O}/\ce{HO-} & eau & ion hydroxyde\\
\ce{CO2,H2O}/\ce{HCO3-} & dioxyde de carbone dissous & ion hydrogénocarbonate\\
\ce{HCO3-}/\ce{CO3^2-} & ion hydrogénocarbonate & ion carbonate\\
\ce{R-COOH}/\ce{R-COO-} & acide carboxylique (ex. : acide éthanoïque) & ion carboxylate (ex. : ion éthanoate)\\
\ce{R-NH3+}/\ce{R-NH2} & ion ammonium (ex. : ion méthylammonium) & amine (ex. : méthylamine)\\
\ce{NH4+}/\ce{NH3} & ion ammonium & ammoniac\\\bottomrule
\end{tabularx}\end{center}
\begin{definition}[Espèce amphotère]
Une espèce \textbf{amphotère} est l'acide d'un couple et la base d'un autre couple : selon
l'espèce avec laquelle elle réagit, elle cède ou capte un ion \ce{H+}. L'eau \ce{H2O} et
l'ion hydrogénocarbonate \ce{HCO3-} sont amphotères.
\end{definition}
\schema{01-amphotere}
\begin{attention}
\ce{CO2,H2O} désigne le dioxyde de carbone dissous dans l'eau, parfois noté \ce{H2CO3}
(acide carbonique). Ce n'est pas le dioxyde de carbone gazeux \ce{CO2(g)}, qui se dégage
lorsque la solution en contient trop : \ce{CO2,H2O(aq) -> CO2(g) + H2O(l)}.
\end{attention}
\essSection{4}{Schémas de Lewis des acides carboxyliques et des amines}
\schema{01-lewis}
\begin{propriete}[Sites acide et basique]
Dans le groupe carboxyle \ce{-COOH}, la liaison O--H est polarisée, car l'oxygène est
plus électronégatif que l'hydrogène : c'est cet atome d'hydrogène qui est cédé sous forme
d'ion \ce{H+}. Dans une amine, l'atome d'azote porte un doublet non liant : il peut capter un
ion \ce{H+} en formant une quatrième liaison. La formule semi-développée suffit souvent :
\ce{CH3-COOH}, \ce{CH3-COO-}, \ce{CH3-NH2}, \ce{CH3-NH3+}.
\end{propriete}
\begin{savoirfaire}[Placer la charge d'un ion]
Pour chaque atome, la charge vaut : (nombre d'électrons de valence) $-$ (nombre d'électrons
des doublets non liants) $-$ (nombre de liaisons).
Oxygène de l'ion carboxylate : $6-6-1=-1$. Azote de l'ion ammonium : $5-0-4=+1$.
La somme des charges des atomes est égale à la charge de l'ion.
\end{savoirfaire}
\essSection{5}{Le pH d'une solution aqueuse}
\begin{definition}[pH]
Le pH d'une solution aqueuse diluée est défini par
\[\pH=-\log\left(\frac{\hho}{\cst}\right)\quad\iff\quad\hho=\cst\times10^{-\pH},\]
où \hho\ est la concentration en ion oxonium, en \molL, et $\cst=\qty{1}{mol.L^{-1}}$ la
concentration standard. Le pH est une grandeur sans unité.
\end{definition}
\begin{propriete}[Lire une variation de pH]
Le pH diminue lorsque \hho\ augmente. Si \hho\ est multipliée par 10, le pH diminue d'une
unité ; si elle est multipliée par 2, il diminue de $\log2\approx0{,}3$.
À 25 °C, une solution est acide si $\pH<7$, neutre si $\pH=7$ et basique si $\pH>7$.
\end{propriete}
\schema{01-echelle-ph}
\begin{savoirfaire}[Calculer avec le bon nombre de chiffres]
\[\hho=\qty{2,5e-4}{mol.L^{-1}}\ \Longrightarrow\ \pH=-\log\left(\num{2.5e-4}\right)=3{,}60\]
\[\pH=8{,}20\ \Longrightarrow\ \hho=10^{-8{,}20}\times\cst=\qty{6,3e-9}{mol.L^{-1}}\]
Le nombre de décimales du pH est égal au nombre de chiffres significatifs de \hho.
\end{savoirfaire}
\begin{experience}[Diluer une solution d'acide chlorhydrique]
Le chlorure d'hydrogène réagit totalement avec l'eau : \ce{HCl(g) + H2O(l) -> H3O^+(aq) + Cl^-(aq)}.
Dans une solution d'acide chlorhydrique de concentration $c$ en acide apporté, on s'attend donc à
$\hho=c$ et à $\pH=-\log(c/\cst)$. On le teste par des dilutions successives d'un facteur 10 :
le pH augmente alors d'environ une unité à chaque dilution, tant que la solution n'est ni trop
concentrée, ni trop diluée. Diluer une solution acide fait tendre son pH vers 7, sans jamais le dépasser.
\schema{01-dilution}
\vspace{1.5mm}
Porter des lunettes et des gants ; étalonner le pH-mètre avec deux solutions étalons ; rincer
la sonde à l'eau distillée entre deux mesures.
\end{experience}
\begin{rediger}
Écrire « l'acide éthanoïque cède un ion hydrogène à l'ammoniac » plutôt que « l'acide donne un
hydrogène ». Distinguer la concentration $c$ en acide apporté et la concentration \hho\ en ion
oxonium. Appuyer chaque conclusion sur une observation ou une donnée : « la solution devient
jaune, donc son pH est inférieur à 6,0 ».
\end{rediger}
\exercices

\rubrique{Découverte}
\begin{blocLivre}{ex-1-1}
\begin{exercice}[1]{Acide ou base ?}
On considère les espèces chimiques suivantes :
\ce{CH3COOH}, \ce{NH3}, \ce{HCOO-}, \ce{CH3NH3+}, \ce{HO-}, \ce{H3O+} et \ce{HCO3-}.
\begin{questions}
\item Pour chaque espèce, indiquer s'il s'agit d'un acide ou d'une base au sens de Brönsted,
puis donner la formule de l'espèce conjuguée.
\item Écrire chaque couple sous la forme acide/base. Quelle espèce de la liste appartient à deux couples ?
\end{questions}
\end{exercice}
\end{blocLivre}
\begin{blocLivre}{ex-1-2}
\begin{exercice}[1]{Écrire des demi-équations}
\begin{questions}
\item Écrire la demi-équation acide-base associée à chacun des couples suivants :
\ce{HCOOH}/\ce{HCOO-} ; \ce{NH4+}/\ce{NH3} ; \ce{H2O}/\ce{HO-} ; \ce{CO2,H2O}/\ce{HCO3-}.
\item Vérifier, pour chacune, la conservation des éléments et de la charge électrique.
\item L'ion hydrogénosulfate \ce{HSO4-} et l'ion sulfate \ce{SO4^2-} forment-ils un couple acide-base ?
Même question pour l'acide sulfurique \ce{H2SO4} et l'ion sulfate.
\end{questions}
\end{exercice}
\end{blocLivre}
\begin{blocLivre}{ex-1-3}
\begin{exercice}[1]{Du pH à la concentration, et inversement}
\begin{questions}
\item Calculer le pH des solutions dont la concentration en ion oxonium vaut :
\begin{tasks}[label=\alph*)](3)
\task \qty{1,0e-3}{mol.L^{-1}} ;
\task \qty{2,5e-4}{mol.L^{-1}} ;
\task \qty{6,3e-9}{mol.L^{-1}}.
\end{tasks}
\item Calculer la concentration en ion oxonium des solutions dont le pH vaut :
\begin{tasks}[label=\alph*)](3)
\task $2{,}0$ ;
\task $7{,}40$ ;
\task $11{,}3$.
\end{tasks}
\item À 25 °C, indiquer pour chacune des six solutions si elle est acide, neutre ou basique.
\end{questions}
\end{exercice}
\end{blocLivre}
\begin{blocLivre}{ex-1-4}
\begin{exercice}[1]{L'acide des fourmis}
L'acide méthanoïque \ce{HCOOH}, sécrété par certaines fourmis, est l'acide carboxylique le plus simple.
\begin{questions}
\item Représenter le schéma de Lewis de la molécule d'acide méthanoïque et entourer son groupe caractéristique.
\item Quel atome d'hydrogène est cédé lors d'une réaction acide-base ? Justifier à l'aide des électronégativités.
\item Représenter le schéma de Lewis de l'ion méthanoate \ce{HCOO-} et justifier la position de sa charge.
\end{questions}
\begin{donnees}
Électronégativités : $\chi(\mathrm{H})=2{,}20$ ; $\chi(\mathrm{C})=2{,}55$ ; $\chi(\mathrm{O})=3{,}44$.
Électrons de valence : H : 1 ; C : 4 ; O : 6.
\end{donnees}
\end{exercice}
\end{blocLivre}
\begin{blocLivre}{ex-1-5}
\begin{exercice}[1]{Espèces amphotères}
\begin{questions}
\item Écrire les deux couples acide-base auxquels appartient l'eau. Justifier qu'elle est amphotère.
\item Écrire l'équation de la réaction entre l'ion oxonium et l'ion hydroxyde.
\item On donne les couples \ce{H3PO4}/\ce{H2PO4-}, \ce{H2PO4-}/\ce{HPO4^2-} et \ce{HPO4^2-}/\ce{PO4^3-}.
Identifier les espèces amphotères.
\item L'ion carbonate \ce{CO3^2-} est-il amphotère ? Justifier à l'aide des couples du cours.
\end{questions}
\end{exercice}
\end{blocLivre}
\rubrique{Entraînement}
\begin{blocLivre}{ex-1-6}
\begin{exercice}[2]{Établir l'équation d'une réaction}
Pour chaque transformation, identifier l'acide et la base qui réagissent, préciser les couples
mis en jeu, puis établir l'équation de la réaction à l'aide des demi-équations.
Les transformations sont considérées comme totales.
\begin{questions}
\item L'acide éthanoïque \ce{CH3COOH} réagit avec l'ammoniac \ce{NH3}.
\item La méthylamine \ce{CH3NH2} réagit avec l'ion oxonium \ce{H3O+}.
\item L'ion hydrogénocarbonate \ce{HCO3-} réagit avec l'ion hydroxyde \ce{HO-}.
\item L'acide méthanoïque \ce{HCOOH} réagit avec l'ion carbonate \ce{CO3^2-}.
\item L'eau est formée dans les réactions 2 et 3. Joue-t-elle le même rôle dans les deux couples concernés ?
\end{questions}
\end{exercice}
\end{blocLivre}
\begin{blocLivre}{ex-1-7}
\begin{exercice}[2]{Souffler dans une solution de bleu de bromothymol}
Un élève souffle avec une paille dans de l'eau additionnée de quelques gouttes de bleu de
bromothymol (BBT). La solution, initialement verte, devient jaune. Il ajoute ensuite quelques
gouttes d'une solution d'hydroxyde de sodium (\ce{Na+}, \ce{HO-}) : la solution devient bleue.
\schema{01-bbt}
\begin{donnees}
Teinte du BBT : jaune si $\pH<6{,}0$ ; verte si $6{,}0<\pH<7{,}6$ ; bleue si $\pH>7{,}6$.
L'air expiré contient environ 4 \% de dioxyde de carbone, contre 0,04 \% dans l'air inspiré.
\end{donnees}
\begin{questions}
\item Comment évoluent le pH, puis la concentration en ion oxonium, lorsque l'élève souffle ?
\item Le dioxyde de carbone expiré se dissout dans l'eau. Écrire l'équation de la réaction,
non totale, entre \ce{CO2,H2O} et l'eau. Préciser les couples et expliquer le changement de teinte.
\item Expliquer le virage au bleu. Écrire l'équation de la réaction entre \ce{CO2,H2O} et
l'ion hydroxyde, considérée comme totale.
\item Quelles espèces amphotères interviennent dans ces transformations ?
\end{questions}
\end{exercice}
\end{blocLivre}
\begin{blocLivre}{ex-1-8}
\begin{exercice}[2]{Une amine et son ion ammonium}
L'éthylamine \ce{CH3-CH2-NH2} est utilisée dans l'industrie des colorants et des médicaments.
\begin{questions}
\item Représenter le schéma de Lewis de l'éthylamine. Repérer le site responsable de son caractère basique.
\item Donner la formule semi-développée et le schéma de Lewis de l'ion éthylammonium, acide conjugué
de l'éthylamine. Justifier la charge.
\item L'éthylamine réagit avec l'eau ; la transformation n'est pas totale. Écrire l'équation de la réaction.
\item Une solution d'éthylamine a un pH de 11,8. Calculer \hho. Ce résultat est-il cohérent avec la question 3 ?
\end{questions}
\end{exercice}
\end{blocLivre}
\begin{blocLivre}{ex-1-9}
\begin{exercice}[2]{L'acide lactique}
L'acide lactique se forme dans les muscles lors d'un effort intense ; il contribue aussi à
l'acidité des yaourts. Sa formule semi-développée est \ce{CH3-CH(OH)-COOH}.
\begin{questions}
\item Recopier la formule, entourer et nommer les groupes caractéristiques.
\item Écrire le couple acide-base de l'acide lactique et la formule semi-développée de l'ion lactate.
\item Représenter le schéma de Lewis de l'ion lactate.
\item Dans le sang, l'acide lactique réagit avec les ions hydrogénocarbonate. Établir l'équation
de la réaction, considérée comme totale. Que devient l'espèce \ce{CO2,H2O} formée ?
\end{questions}
\end{exercice}
\end{blocLivre}
\begin{blocLivre}{ex-1-10}
\begin{exercice}[2]{Une eau minérale riche en ions hydrogénocarbonate}
Une eau minérale gazeuse porte l'indication : « hydrogénocarbonates \ce{HCO3-} : 3\,050 mg/L ».
Cette eau est parfois conseillée en cas d'acidité gastrique : le suc gastrique contient de l'acide
chlorhydrique, donc des ions \ce{H3O+}.
\begin{donnees}
$M(\mathrm{H})=\qty{1,0}{g.mol^{-1}}$ ; $M(\mathrm{C})=\qty{12,0}{g.mol^{-1}}$ ; $M(\mathrm{O})=\qty{16,0}{g.mol^{-1}}$ ;
volume molaire des gaz à 20 °C sous 1 bar : $V_\mathrm{m}=\qty{24,0}{L.mol^{-1}}$.
\end{donnees}
\begin{questions}
\item Montrer que la concentration en ion hydrogénocarbonate de cette eau vaut \qty{5,00e-2}{mol.L^{-1}}.
\item Justifier que l'ion hydrogénocarbonate est une espèce amphotère.
\item Écrire l'équation de la réaction, considérée comme totale, entre les ions \ce{HCO3-} et \ce{H3O+}.
Pourquoi cette réaction peut-elle provoquer des éructations ?
\item Calculer la quantité maximale d'ions \ce{H3O+} consommée par un verre de \qty{200}{mL} de cette eau,
puis le volume maximal de gaz qui peut se dégager.
\end{questions}
\end{exercice}
\end{blocLivre}
\begin{blocLivre}{ex-1-11}
\begin{exercice}[2]{Lire une variation de pH}
\begin{donnees}
$\log2\approx0{,}30$ ; $\log5\approx0{,}70$ ; $10^{0{,}1}\approx1{,}26$ ; $10^{0{,}3}\approx2{,}0$ ; $10^{0{,}5}\approx3{,}2$.
\end{donnees}
\begin{questions}
\item Le pH d'une eau de pluie « naturelle » est voisin de 5,6. Dans une région industrielle, on mesure
un pH de 4,3. Par quel facteur la concentration en ion oxonium a-t-elle été multipliée ?
\item Le pH du sang est normalement compris entre 7,35 et 7,45. Calculer le rapport des concentrations
en ion oxonium correspondant à ces deux valeurs. Commenter.
\item Une solution d'acide chlorhydrique a un pH de 2,7. On la dilue cinq fois. En admettant que \hho\
est divisée par 5, prévoir le pH de la solution diluée.
\item Par quel facteur faut-il diluer une solution d'acide chlorhydrique pour augmenter son pH de 0,5 ?
\end{questions}
\end{exercice}
\end{blocLivre}
\rubrique{Objectif bac}
\begin{blocLivre}{ex-1-12}
\begin{exercice}[3]{L'odeur du poisson}
\begin{docu}{1}{Une amine volatile}
L'odeur du poisson qui n'est plus frais est due principalement à la triméthylamine \ce{(CH3)3N},
une amine volatile formée lors de la dégradation bactérienne de composés présents dans la chair du poisson.
\end{docu}
\begin{docu}{2}{Un conseil de cuisinier}
Pour atténuer l'odeur, on arrose le poisson de jus de citron ou de vinaigre. Le vinaigre contient
de l'acide éthanoïque \ce{CH3COOH}. Les ions, contrairement à la triméthylamine, ne passent
pratiquement pas dans l'air : ils restent dans la phase aqueuse.
\end{docu}
\begin{questions}
\item Représenter le schéma de Lewis de la triméthylamine. Quel est le site basique de la molécule ?
\item Écrire le couple acide-base de la triméthylamine. Représenter le schéma de Lewis de l'acide conjugué
et justifier sa charge.
\item Établir l'équation de la réaction, considérée comme totale, entre la triméthylamine et l'acide éthanoïque.
\item Expliquer pourquoi le vinaigre atténue l'odeur du poisson.
\item Le jus de citron contient de l'acide citrique, noté AH. Écrire l'équation de sa réaction avec la
triméthylamine. Pourquoi de l'eau salée (\ce{Na+}, \ce{Cl-}) n'aurait-elle pas le même effet ?
\end{questions}
\end{exercice}
\end{blocLivre}
\begin{blocLivre}{ex-1-13}
\begin{exercice}[3]{Un ballon qui se gonfle tout seul}
Dans une bouteille, on verse $V=\qty{50,0}{mL}$ de vinaigre blanc à 6,0 \%. On place $m=\qty{5,00}{g}$
d'hydrogénocarbonate de sodium \ce{NaHCO3(s)} dans un ballon de baudruche que l'on fixe sur le goulot.
On redresse le ballon : la poudre tombe dans le vinaigre, une effervescence se produit et le ballon se gonfle.
\schema{01-ballon}
\begin{donnees}
Vinaigre à 6,0 \% : 6,0 g d'acide éthanoïque \ce{CH3COOH} pour 100 g de vinaigre ; masse volumique du vinaigre :
$\rho=\qty{1,01}{g.mL^{-1}}$ ; $M(\ce{CH3COOH})=\qty{60,0}{g.mol^{-1}}$ ; $M(\ce{NaHCO3})=\qty{84,0}{g.mol^{-1}}$ ;
$V_\mathrm{m}=\qty{24,0}{L.mol^{-1}}$ ; dissolution : \ce{NaHCO3(s) -> Na^+(aq) + HCO3^-(aq)} ;
volume d'une boule de rayon $R$ : $\frac43\pi R^3$.
\end{donnees}
\begin{questions}
\item Établir l'équation de la réaction, considérée comme totale, entre l'acide éthanoïque et l'ion
hydrogénocarbonate. Quel gaz provoque l'effervescence ?
\item Montrer que la concentration en acide éthanoïque apporté dans le vinaigre vaut environ \qty{1,0}{mol.L^{-1}}.
\item Déterminer le réactif limitant et l'avancement maximal.
\item \initiative Le ballon peut-il atteindre un diamètre de 15 cm ? On admet que tout le dioxyde de carbone
formé passe dans le ballon, assimilé à une boule.
\item Proposer des quantités de vinaigre et d'hydrogénocarbonate de sodium permettant d'atteindre ce diamètre.
\end{questions}
\end{exercice}
\end{blocLivre}
\begin{blocLivre}{ex-1-14}
\begin{exercice}[3]{L'acidification des océans}
\begin{docu}{1}{Un pH qui diminue}
Le pH moyen de l'eau de mer de surface est passé de 8,11 en 1985 à 8,04 en 2024. Selon le scénario
d'émissions de dioxyde de carbone, la baisse projetée d'ici 2100 est comprise entre 0,15 et 0,5 unité pH.
\source{D'après l'Agence européenne pour l'environnement, indicateur « Ocean acidification », 25 août 2026, données du Copernicus Marine Service.}
\end{docu}
\begin{docu}{2}{Des coquilles en carbonate de calcium}
Le dioxyde de carbone atmosphérique se dissout dans l'océan. De nombreux organismes marins (coraux,
mollusques, plancton calcaire) fabriquent leur squelette ou leur coquille à partir des ions carbonate :
\ce{Ca^2+(aq) + CO3^2-(aq) -> CaCO3(s)}.
\end{docu}
\schema{01-ocean}
\begin{questions}
\item Écrire les couples acide-base dont fait partie l'ion hydrogénocarbonate. Quelle propriété possède-t-il ?
\item Écrire l'équation de la réaction, non totale, entre le dioxyde de carbone dissous et l'eau.
Expliquer pourquoi la dissolution du dioxyde de carbone fait baisser le pH.
\item Le dioxyde de carbone dissous réagit aussi avec les ions carbonate. Établir l'équation de cette
réaction. Pourquoi l'augmentation de la teneur en \ce{CO2} de l'atmosphère menace-t-elle les organismes à coquille ?
\item Calculer \hho\ en 1985 et en 2024. Vérifier que l'augmentation relative est voisine de 18 \%.
\item En admettant que la baisse projetée s'ajoute à la valeur de 2024, encadrer \hho\ en 2100
et le facteur d'augmentation par rapport à 2024.
\item Un article affirme : « D'ici 2100, les océans vont devenir acides. » Commenter cette phrase.
\end{questions}
\end{exercice}
\end{blocLivre}
\begin{blocLivre}{ex-1-15}
\begin{exercice}[3]{Tester la relation entre pH et concentration}
On cherche à tester la relation $\pH=-\log(c/\cst)$ pour des solutions d'acide chlorhydrique de
concentration $c$ en acide apporté, préparées par dilutions successives d'un facteur 10 à partir d'une
solution S$_0$ de concentration $c_0=\qty{1,0e-1}{mol.L^{-1}}$.
\begin{donnees}
Matériel : pipettes jaugées de 5,0 mL, 10,0 mL et 20,0 mL ; fioles jaugées de 50,0 mL, 100,0 mL et
250,0 mL ; béchers ; pH-mètre et solutions étalons de pH 4,00 et 7,00 ; eau distillée.
\end{donnees}
\begin{questions}
\item \capacite{Expérience}Rédiger le protocole de préparation de 100,0 mL de solution S$_1$ de concentration
$c_0/10$. Justifier le choix de la verrerie.
\item Citer trois précautions à prendre lors des mesures de pH.
\item \capacite{Capacité numérique}On obtient les mesures ci-dessous. Le programme réalise un ajustement
affine $\pH=a\times(-\log(c/\cst))+b$ sur les cinq premières mesures.
\begin{center}\small\renewcommand{\arraystretch}{1.2}
\begin{tabular}{@{}l*{6}{c}@{}}\toprule
$c$ (\molL) & \num{1.0e-1} & \num{1.0e-2} & \num{1.0e-3} & \num{1.0e-4} & \num{1.0e-5} & \num{1.0e-6}\\
pH mesuré & 1,10 & 2,04 & 3,02 & 4,01 & 4,98 & 5,81\\\bottomrule
\end{tabular}\end{center}
\begin{codecote}[0.66]
\begin{pythoncol}
import numpy as np
c = np.array([1e-1, 1e-2, 1e-3, 1e-4, 1e-5])
pH = np.array([1.10, 2.04, 3.02, 4.01, 4.98])
x = -np.log10(c)
a, b = np.polyfit(x, pH, 1)
print(round(a, 3), round(b, 3))
\end{pythoncol}
\cote
\begin{consolecol}
0.973 0.111
\end{consolecol}
\end{codecote}
Ces résultats sont-ils compatibles avec la relation testée ?
\item Proposer une explication à l'écart observé pour la solution la plus concentrée, puis pour la plus diluée.
\item \initiative Dans quel domaine de concentrations la relation est-elle vérifiée à 0,05 unité pH près ?
\end{questions}
\end{exercice}
\end{blocLivre}
